% GENERATED FILE. Edit canonical lessons, then run npm run generate:books.
% canonical-source-hash: fnv1a64:e5779f1ea3d9ee68
% canonical-lessons: RU-C16-i, RU-C16-i-predlozheniya, RU-C16-ili, RU-C16-ni-ni, RU-C16-tozhe

\chapter{And, Or, and Neither}
\label{ch:ru-and-or-neither}
\hlchaptermodality{16}

\begin{chapteropening}
\textbf{By the end of this chapter:} \emph{I can join two words or two whole sentences with i, offer a choice with ili, refuse both halves at once with ni ... ni, and agree with somebody using tozhe.}

The first two-clause sentence this book can make. Until this chapter the longest structure the track taught was a single clause, and i -- the commonest word in Russian -- had been printed twice as a conjunction without any lesson owning it. ni ... ni is the first joining pattern in the book that stands in FRONT of both halves rather than between them, and it needs the negative said three times in one short sentence.
\end{chapteropening}

\section[i]{\ru{и} — the word you have already read}

\label{lesson:RU-C16-i}



\begin{quote}
Two verbs from the chapter of first verbs, before the smallest word in the book arrives.

\begin{itemize}
\raggedright
  \item \emph{Recall:} \emph{\ru{знать}}, its English cousin \emph{know}, and how \textbf{\ru{не}} goes in front
  \item \emph{Recall:} \emph{\ru{говорить}}, which is \textbf{not} a relative of \emph{govern}, and the two verb families it belongs to
  \item \emph{Recall:} the eight letters you sorted into true friends, false friends and new shapes
\end{itemize}
\end{quote}

\subsection*{You'll want to know first — \ru{и}}
\begin{itemize}
\raggedright
  \item \textbf{\ru{и}} — \emph{i} — \textbf{and}
\end{itemize}

One letter. One vowel. Nothing else. It is the commonest word in Russian, and you have been reading it since the verbs of the mind:

\begin{itemize}
\raggedright
  \item \textbf{\ru{Я} \ru{читаю} \ru{и} \ru{понимаю}.} — \emph{ya chitáyu i panimáyu} — I read \textbf{and} I understand.
\end{itemize}

That sentence was on the page long before anyone told you that the little \textbf{\ru{и}} between the two verbs was a word at all. It is. It joins two things of the same kind: two nouns, two verbs, two whole thoughts.

The letter is one you can already write, and the word is the letter. Nothing new goes on the page here — only the fact that this shape has a meaning.

\begin{cousinweb}
\textbf{\ru{и}} goes back to a Slavic particle built on the ancient \textbf{pointing stem} \emph{*i-}, the demonstrative that meant roughly \emph{this one, that one}. A joining word that began as a pointing word is not strange: \emph{and} is a way of saying \emph{there is another one over here}.

That same old stem is the one hiding in \textbf{\ru{его}} — which is why the smallest word in the language and one of its pronouns are relatives.
\end{cousinweb}

\subsection*{Your turn}
\begin{itemize}
\raggedright
  \item \emph{Say it:} \emph{\ru{я} \ru{читаю} \ru{и} \ru{понимаю}} — and hear \textbf{\ru{и}} as a word this time
  \item \emph{Say it:} any two verbs you own, with \textbf{\ru{и}} between them
  \item \emph{Build:} two foods, joined — \emph{\ru{молоко} \ru{и} \ru{сок}}
\end{itemize}

\begin{tcolorbox}[breakable,colback=teal!4,colframe=teal!35!black,title={Before you move on}]
Say \textquotedblleft{}and.\textquotedblright{} (\textbf{\ru{и}}.) Which letter is it, and had you written that letter before today? (\textbf{\ru{и}} — yes, from the writing lessons.) What kind of old word is it built on? (\textbf{A pointing word.}) And one look back: say \textquotedblleft{}I know,\textquotedblright{} and then say you do not. (\textbf{\ru{я} \ru{знаю}} / \textbf{\ru{я} \ru{не} \ru{знаю}}.)
\end{tcolorbox}
\section[… i …]{\ru{и} between two sentences — the book's first long sentence}

\label{lesson:RU-C16-i-predlozheniya}



\begin{quote}
Yesterday's word, and two verbs it is about to join.

\begin{itemize}
\raggedright
  \item \emph{Recall:} \emph{\ru{идти}}, the one-way verb whose past comes from another root entirely
  \item \emph{Recall:} \emph{\ru{думать}}, the verb hiding inside the name of the Russian parliament, and the partner verb that finishes what it starts
\end{itemize}
\end{quote}

\begin{grammarlens}[title={Grammar lens — one word, any size of piece}]
You have joined two verbs with \textbf{\ru{и}}. Now join two whole sentences with it. The word does not change:

\begin{itemize}
\raggedright
  \item \textbf{\ru{я} \ru{иду} \ru{и} \ru{я} \ru{думаю}} — \emph{ya idú i ya dúmayu} — I am walking \textbf{and} I am thinking.
\end{itemize}

Each half stands on its own. Take the \textbf{\ru{и}} away and you have two sentences that are still Russian. Put it back and you have one sentence with two halves, which is a thing this book has not been able to make until now.

Russian does not need a comma here, and does not want one: when \textbf{\ru{и}} joins two halves that share a subject, nothing goes in front of it.

\begin{itemize}
\raggedright
  \item \textbf{\ru{я} \ru{иду} \ru{и} \ru{думаю}} — \emph{ya idú i dúmayu} — I walk and think.
\end{itemize}

Say the second \textbf{\ru{я}} or leave it out; both are Russian. Leaving it out is what a Russian usually does, because the ending on the verb has already said who.
\end{grammarlens}

\subsection*{Your turn}
\begin{itemize}
\raggedright
  \item \emph{Say it:} \emph{\ru{я} \ru{иду} \ru{и} \ru{думаю}}
  \item \emph{Build:} any two things you do, joined with \textbf{\ru{и}}, dropping the second \emph{\ru{я}}
  \item \emph{Build:} the same pair with the second \emph{\ru{я}} put back
\end{itemize}

\begin{tcolorbox}[breakable,colback=teal!4,colframe=teal!35!black,title={Before you move on}]
Does \textbf{\ru{и}} change when the pieces it joins get bigger? (\textbf{No} — one letter, always, and it is still the same old pointing word underneath.) Does a comma go in front of it when both halves share a subject? (\textbf{No}.) And one look back: which verb of motion means going one way, this once? (The verb is \textbf{\ru{идти}}.)
\end{tcolorbox}
\section[íli]{\ru{или} — \textquotedblleft{}or,\textquotedblright{} which is \textquotedblleft{}and\textquotedblright{} with a doubt on the end}

\label{lesson:RU-C16-ili}



\begin{quote}
The two verbs that first stood on either side of \textbf{\ru{и}}.

\begin{itemize}
\raggedright
  \item \emph{Recall:} \emph{\ru{понимать}}, which literally means to \textbf{take hold of} — the same picture English keeps in \emph{comprehend}
  \item \emph{Recall:} \emph{\ru{читать}}, and the fact that Russian has \textbf{no separate way} to say \emph{I am reading}
  \item \emph{Recall:} the whole first exchange — \emph{\ru{да}}, \emph{\ru{нет}}, and which greeting you give a stranger
\end{itemize}
\end{quote}

\subsection*{You'll want to know first — \ru{или}}
\begin{itemize}
\raggedright
  \item \textbf{\ru{или}} — \emph{íli} — \textbf{or}
\end{itemize}

Stress the first syllable: \emph{Í-li}. It sits between the two things exactly where \textbf{\ru{и}} sat, and it does the opposite job — not both, but one of the two.

\begin{itemize}
\raggedright
  \item \textbf{\ru{чай} \ru{или} \ru{кофе}?} — \emph{chai íli kófe} — tea \textbf{or} coffee?
  \item \textbf{\ru{я} \ru{читаю} \ru{или} \ru{пишу}} — \emph{ya chitáyu íli pishú} — I read \textbf{or} I write.
\end{itemize}

Offering somebody a choice is one of the first useful things you can do with a joining word, and this is the whole of it: two nouns, \textbf{\ru{или}}, a rising voice.

\begin{cousinweb}
\textbf{\ru{или}} is \textbf{\ru{и}} with a small word welded on the end: \textbf{\ru{ли}}, an old particle that carries doubt — the thing Russian puts in a sentence when the answer is not yet known.

So \emph{and} plus \emph{maybe} gives \emph{or}. Once you see \textbf{\ru{и}} at the front of \textbf{\ru{или}}, the pair is impossible to confuse: the short one joins, the long one doubts.
\end{cousinweb}

\subsection*{Your turn}
\begin{itemize}
\raggedright
  \item \emph{Say it:} \emph{\ru{чай} \ru{или} \ru{кофе}?} with the voice going up
  \item \emph{Build:} the same two drinks joined with \textbf{\ru{и}} instead — and hear the meaning turn over
  \item \emph{Build:} two of your own verbs with \textbf{\ru{или}} between them
\end{itemize}

\begin{tcolorbox}[breakable,colback=teal!4,colframe=teal!35!black,title={Before you move on}]
Say \textquotedblleft{}or.\textquotedblright{} (It is \textbf{\ru{или}}.) Which two pieces is it built from? (\textbf{\ru{и}} plus \textbf{\ru{ли}}.) Which of the two, \textbf{\ru{и}} or \textbf{\ru{или}}, offers a choice? (The long one, \textbf{\ru{или}}.) And one look back: what does \emph{\ru{понимать}} literally do to a thing? (\textbf{Takes hold of it.})
\end{tcolorbox}
\section[ni … ni …]{\ru{ни} … \ru{ни} … — refusing both halves at once}

\label{lesson:RU-C16-ni-ni}



\begin{quote}
Two verbs, and the letter one of them brought with it.

\begin{itemize}
\raggedright
  \item \emph{Recall:} \emph{\ru{писать}}, whose root means to \textbf{scratch} — the same root behind \emph{paint} and \emph{picture} — and the letter \textbf{\ru{ш}} it taught you
  \item \emph{Recall:} \emph{\ru{брать}}, English \emph{bear}, and the partner verb from a completely different root
  \item \emph{Recall:} \emph{\ru{или}}, and the two pieces \textbf{\ru{и}} and \textbf{\ru{ли}} it is welded from
\end{itemize}
\end{quote}

\begin{grammarlens}[title={Grammar lens — the word that goes in front, twice}]
\textbf{\ru{и}} and \textbf{\ru{или}} both stand \textbf{between} the two things. This one stands \textbf{in front} of each of them, and there are two of it:

\begin{itemize}
\raggedright
  \item \textbf{\ru{ни} \ru{чай}, \ru{ни} \ru{кофе}} — \emph{ni chai, ni kófe} — \textbf{neither} tea \textbf{nor} coffee.
\end{itemize}

That shape is new. Every joining word so far has been a hinge in the middle; \textbf{\ru{ни} … \ru{ни} …} is a pair of flags planted before each half.

\textbf{\ru{ни}} is the negator \textbf{\ru{не}} with its vowel changed. Russian keeps them apart by job: \textbf{\ru{не}} denies one thing, \textbf{\ru{ни}} is the form that turns up when the denial has to be said twice.

And Russian says the denial a third time, on the verb. Where English is content with \emph{I take neither}, Russian negates the verb as well:

\begin{itemize}
\raggedright
  \item \textbf{\ru{я} \ru{не} \ru{беру} \ru{ни} \ru{чай}, \ru{ни} \ru{сок}} — \emph{ya ne berú ni chai, ni sok} — I take neither tea nor juice.
\end{itemize}

Three negatives in one short sentence, and every one of them is required. That is not sloppiness — it is the rule.
\end{grammarlens}

\subsection*{Your turn}
\begin{itemize}
\raggedright
  \item \emph{Say it:} \emph{\ru{ни} \ru{чай}, \ru{ни} \ru{кофе}}
  \item \emph{Say it:} the whole refusal — \emph{\ru{я} \ru{не} \ru{беру} \ru{ни} \ru{чай}, \ru{ни} \ru{сок}}
  \item \emph{Build:} two foods you own, refused together
\end{itemize}

\begin{tcolorbox}[breakable,colback=teal!4,colframe=teal!35!black,title={Before you move on}]
Where does \textbf{\ru{ни}} stand — between the halves or in front of them? (\textbf{In front of each half}.) How many times is the negative said in \emph{\ru{я} \ru{не} \ru{беру} \ru{ни} \ru{чай}, \ru{ни} \ru{сок}}? (\textbf{Three}.) And one look back: which verb means to write, and which letter arrived with it? (The verb is \textbf{\ru{писать}}; the letter is \textbf{\ru{ш}}.)
\end{tcolorbox}
\section[tózhe]{\ru{тоже} — \textquotedblleft{}me too,\textquotedblright{} and \textquotedblleft{}me neither\textquotedblright{}}

\label{lesson:RU-C16-tozhe}



\begin{quote}
Two verbs whose roots went a long way from home.

\begin{itemize}
\raggedright
  \item \emph{Recall:} \emph{\ru{спрашивать}}, on the root that gave Latin its word for \textbf{pray}
  \item \emph{Recall:} \emph{\ru{помогать}}, built on a verb meaning \textbf{to be able} — English \emph{may}, \emph{might} and \emph{main}
\end{itemize}
\end{quote}

\subsection*{You'll want to know first — \ru{тоже}}
\begin{itemize}
\raggedright
  \item \textbf{\ru{тоже}} — \emph{tózhe} — \textbf{also, too}
\end{itemize}

Stress the first syllable. It is not a joining word — it does not sit between two things — but it does the same social work as one: it puts a second person or a second fact alongside the first.

\begin{itemize}
\raggedright
  \item \textbf{\ru{я} \ru{тоже}} — \emph{ya tózhe} — \textbf{me too}.
\end{itemize}

Two syllables, and one of the most useful answers in the language. Somebody says what they do; you say \textbf{\ru{я} \ru{тоже}} and you have agreed, in full, without owning a single extra verb.

\begin{cousinweb}
\textbf{\ru{тоже}} is two words fused: \textbf{\ru{то}} (that) and \textbf{\ru{же}} (a little word meaning \emph{the very same}). \emph{That same} — which is what \emph{also} has always meant.

And because Russian negates by putting \textbf{\ru{не}} in front of the verb, \emph{me neither} comes out by putting \textbf{\ru{не}} after \textbf{\ru{тоже}}:

\begin{itemize}
\raggedright
  \item \textbf{\ru{я} \ru{тоже} \ru{не} \ru{знаю}} — \emph{ya tózhe ne znáyu} — I do \textbf{not} know \textbf{either}.
\end{itemize}

English swaps the whole word — \emph{too} becomes \emph{either}. Russian leaves \textbf{\ru{тоже}} where it is and lets \textbf{\ru{не}} do the work, which is one fewer word to learn.
\end{cousinweb}

\subsection*{Your turn}
\begin{itemize}
\raggedright
  \item \emph{Say it:} \emph{\ru{я} \ru{тоже}}
  \item \emph{Say it:} \emph{\ru{я} \ru{тоже} \ru{не} \ru{знаю}}
  \item \emph{Build:} agree with any sentence you can already make, using \textbf{\ru{тоже}}
\end{itemize}

\begin{tcolorbox}[breakable,colback=teal!4,colframe=teal!35!black,title={Before you move on}]
Say \textquotedblleft{}me too.\textquotedblright{} (That is \textbf{\ru{я} \ru{тоже}}.) Which two words is \textbf{\ru{тоже}} built from? (\textbf{\ru{то}} and \textbf{\ru{же}}.) How do you turn it into \textquotedblleft{}me neither\textquotedblright{}? (\textbf{Put \ru{не} in front of the verb} — the word \textbf{\ru{тоже}} itself does not change.) And one look back: which verb means to help, and what is buried inside it? (The verb is \textbf{\ru{помогать}}, and the piece inside it means \textbf{to be able}.)
\end{tcolorbox}
